\section{Related Work}
\label{sec:related}

\subsection{Scheduling with predictions}

Purohit, Svitkina and Kumar \cite{purohit2018predictions} frame predictions in online algorithms as a consistency--robustness trade-off. Predicted-size scheduling has since been studied through stochastic classes \cite{mitzenmacher2020misprediction}, bounded multiplicative error \cite{scully2022uniform}, worst-case distortion \cite{azar2021flowtime,azar2022distortion}, many-server limits \cite{dong2025sjf}, prediction acquisition \cite{shahout2024skippredict}, and per-job stretch on a preemptive machine \cite{zhao2025fairpredictions}; the survey of Mitzenmacher and Shahout \cite{mitzenmacher2025queueing} synthesises this line. Our reference is FCFS rather than an optimum, our scope is one job rather than a mean, arrivals and prediction errors are arbitrary, and service is non-preemptive. The resulting bound is correspondingly limited: it controls additional waiting relative to FCFS and does not assert an improvement. Its usefulness is empirical and depends jointly on size dispersion, predictability and arrival structure, not on membership in a necessary heavy-tail class.

\subsection{Fairness relative to FCFS, and bounded overtaking}

Worst-case per-job fairness predates fair-start-time metrics. Schwiegelshohn and Yahyapour \cite{schwiegelshohn2000fairness} call a schedule $\lambda$-fair when jobs submitted after job $i$ cannot increase $i$'s flow time by more than a factor $\lambda$, and bound fairness and efficiency together for a preemptive gang policy. Their bound is multiplicative and, for a short job, tighter than an $L$-denominated one; it is obtained by gang preemption. Our guarantee is instead additive, non-preemptive, and referenced to the same input under FCFS. Proposition~\ref{prop:floor} says only that a wrapper of our family, protecting against every opaque base proposal sequence, cannot uniformly promise below order $L$ once it permits overtaking; it is not an impossibility for all scheduling mechanisms. Supplementary Section~S9 gives their constants and the rest of the per-paper detail this section compresses.
% Source for the constants of this paragraph: they are stated in full in Supplementary
% Section S9, which carries the same source comment.

Fair-start-time work gives a closer reference: a job is compared with a re-simulation that excludes later arrivals \cite{sabin2004fairness}, and the resulting excess is then used as a job-level metric \cite{sabin2005unfairness,rajbhandary2013fairness}. Conservative backfilling instead fixes a reservation for every queued job, so later backfills cannot delay that promised start \cite{lindsay2013backfilling}; it is a $G=0$ endpoint relative to the reservation, not relative to same-input FCFS. PV-EASY \cite{yuan2010pveasy,yuan2014strictfairness} enforces zero lower-priority delay by killing and restarting obstructing backfilled jobs. These guarantees are online but use different references or preemption, whereas ours limits a non-preemptive wrapper's additive FCFS excess.

The closest architectural precedent for a wrapper around an arbitrary policy is Load Balancing Guardrails \cite{grosof2019guardrails}, which augments any dispatching policy so that a load-balancing system with SRPT at each server is mean-response-time optimal in heavy traffic. Its reference, stochastic objective, preemptive server policy and dispatch-level constraints differ from our per-job additive promise, but the separation between an unrestricted base decision and a protective wrapper is shared. At one server our ``pointwise maximal'' statement uses the established supervisory-control meaning of a maximally permissive supervisor, disabling only the actions needed to keep behaviour within the legal specification \cite{ramadge1987supervisory}; we claim that property only for the opaque-base, integer-time class of Theorem~\ref{thm:k1_maximal}, not as a general least-restrictive result. Supplementary Table~S21 compares the reference discipline, arrivals, preemption, information, server model and guarantee across these nearest mechanisms.
% Reading depth for this section is recorded with the extended detail in Supplementary
% Section S9.

Bounded overtaking has a history in queueing theory, in two currencies. Counting came first: Nudge lets a small arriving job swap past the job at the tail of the queue \cite{grosof2021nudge}, Nudge-$K$ lets it pass up to $K$ others while no job is passed more than once \cite{vanhoudt2022nudgek}, the line continues with Nudge*(M) \cite{charlet2026nudgestarm}, and the finite-skip models of WCFS restrict service to the oldest $n$ jobs and add work conservation and non-idling \cite{grosof2022wcfs}. Work is the second currency: $\gamma$-Boost measures the \emph{crossing work} of later arrivals that precede a tagged job \cite{yu2024boost}, and that line extends to $k$ servers \cite{yu2025twotraffics}. All of them are analysed in steady state, under Poisson arrivals with light-tailed or phase-type sizes, and target stationary means or tail constants rather than a deterministic per-job cap enforced on every sample path. That Nudge implies, on every sample path, a per-job additive bound of one small-job length is our own reading and not a claim of \cite{grosof2021nudge}, whose analysis is stochastic by design, and we do not claim bounded overtaking at one server as new.

Within that conjunction, Theorem~\ref{thm:two_size_achievability} supplies a
matching achievable band on a specified batch family, and
Theorem~\ref{thm:k1_maximal} supplies pointwise maximal permission in its
opaque-base integer-time class. Neither claims global instance optimality or a
least-restrictive wrapper for $k\geq2$.

Our contribution is therefore a conjunction and not a new currency. The identity is two-sided and holds on an arbitrary arrival sequence with arbitrary service times, under no distributional assumption; the constant $(2-2/k)L$ is exact for $k$ non-preemptive servers, as is the $(3-2/k)L$ of the wrapper; the converse makes a work budget necessary and not merely sufficient, up to an additive $O(kL)$; and Algorithm~\ref{alg:guard} enforces the bound online around any base policy, learned ones included, under predictions that may be arbitrarily wrong.

The additive shape and the constants are familiar on their own. Two-sided additive bounds against a fluid or ideal reference are the standard form in packet scheduling \cite{bennett1996wf2q,parekh1993gps,parekh1994gps,stoica1995eevdf}, where the error term is one maximum-size packet precisely because that packet cannot be preempted \cite{bennett2002psrg}; the coefficient $2-1/k$ of Graham's list-scheduling bound has the shape of the $(2-2/k)$ and $(3-2/k)$ above, and the additive form appears inside his proof \cite{graham1969bounds}, while non-preemptive real-time feasibility conditions charge a whole uninterrupted execution for the same reason \cite{jeffay1991nonpreemptive}. What we prove is not a conservation law in the sense of Kleinrock \cite{kleinrock1965conservation} and the achievable-region work after it, whose sample-path members are aggregate or class-level identities rather than per-job comparisons against a named reference \cite{green2000samplepath,eltaha2017workload}; the law does extend to several servers \cite{bartsch1978conservation}, but under non-preemptive priority with $k>1$ it needs every class to share one service-time distribution \cite{federgruen1988mgc}, which is exactly what our model withholds. Lemma~\ref{lem:gap} is the quantitative substitute, and it differs from the nearest pathwise work-difference bound \cite[Lemma~2]{grosof2018srpt} in comparing two arbitrary non-preemptive work-conserving $k$-server policies, in bounding total unfinished work, and in being two-sided. Aging disciplines \cite{kleinrock1964delay,stanford2014apq}, the unfairness analyses of SRPT against processor sharing \cite{bansal2001srpt,wierman2003classifying} and the multiserver accumulating priority queue \cite{sharif2014mapq,li2016mapq} describe classes in steady state and not individual jobs, whereas our $\mathrm{over}[q]$ accumulates overtaken work and carries a pathwise cap. The word itself needs one caution: in queueing theory \emph{overtaking} means departing before a customer who arrived earlier \cite{whitt1984overtaking,kim2016overtaking}, while we measure the work that passes a waiting job. A single-server argument also stops working with several servers, the workload vector recursion being Kiefer and Wolfowitz's \cite{kiefer1955queues} and its minimality under first-come first-served Daley's \cite{daley1987fcfs}, which holds only where arrivals are allocated to servers in a manner that does not depend on their service times; our identity replaces that step with a quantitative comparison. Supplementary Section~S9 states each of these bounds in the form its source gives it.

\subsection{Runtime prediction and scheduling in deployed systems}

Backfilling schedulers on high-performance computing systems have used predicted runtimes in place of user estimates for two decades \cite{tsafrir2007backfilling}, with predictions learned from past jobs and queue reordering tuned online \cite{gaussier2015backfilling,gaussier2018easytuning}; their resources are multi-dimensional and their reported metric is a mean bounded slowdown. Closest to our setting, Zhang, Wu and Lu \cite{zhang2024ojscheduling} predict the running time of submitted code in a judging queue and dispatch in ascending predicted runtime inside fixed-size windows, so we are not the first to schedule a judging queue by predicted cost; a fixed window does limit overtaking, but by position, and that work reports no per-job bound and no adversarial-prediction test. In inference serving, where non-preemption is close to a hard constraint, recent systems learn a relative order instead of a length \cite{fu2024ltr} or target the tail because prediction-driven policies are fragile under distribution shift \cite{li2026tailaware,luo2025asrpt}, and production systems otherwise avoid starvation with aging counters and maximum waiting times, without proofs. Serverless platforms are in the same position, SFS approximating shortest-remaining-time-first from user space \cite{fu2022sfs} and ALPS learning a per-function rank and time-slice cap from an offline replay \cite{fu2024alps}, neither with a bound. Our guard is a wrapper rather than a policy, so it can be placed around any of these schedulers, and its bound does not depend on which one is underneath.

\subsection{Time-varying and deadline-driven load}

That people concentrate their actions near a cutoff has been measured directly, in time management research \cite{konig2005deadline} and in conference registration data \cite{alfi2007deadline,alfi2009deadline}. The queueing consequences of an arrival rate that varies on the timescale of the service time are also well studied: peak congestion in a multi-server system lags the peak of the arrival rate \cite{massey1997peak}, and staffing rules that follow a time-varying rate rest on offered-load approximations rather than stationary formulas \cite{jennings1996staffing,green2007coping,feldman2008staffing}. These results say how many servers to provide. They give no waiting bound for an individual job and no guidance for a queue that has already formed, and we leave capacity fixed. The steady-state classified-priority formula that motivates ranking by predicted cost \cite{mitzenmacher2020misprediction,harcholbalter2013} does not survive contact with this load either; Section~\ref{sec:model} measures the mismatch, which is why we state the guarantee on sample paths and assume nothing about arrivals.

\subsection{Relational and tabular models for cost prediction}

Our comparison includes a heterogeneous graph neural network over users, tasks and assignments, in the relational graph convolution family \cite{schlichtkrull2018rgcn,xu2019gin}. Two features of the published benchmarks make them a poor argument either way. In RelBench \cite{robinson2024relbench} and its successor \cite{gu2026relbenchv2} the gradient-boosted baseline is fitted on the raw columns of the entity table alone, without hand-built aggregates over related tables, so the margin reported there is not a comparison against a tabular model carrying engineered history features; and the same paper's study against a human data scientist, which is the closer comparison, records regression as the region where the relational model struggles. Cost prediction is a regression task on entities with long individual histories, which places it in that region, and the evidence that tree ensembles remain strong on tabular data concerns single tables \cite{grinsztajn2022tabular}. We therefore settle the question inside our own protocol instead of importing a conclusion, with a tabular baseline that carries engineered causal history features \cite{ke2017lightgbm}, an ablation ladder that removes and then permutes the edges, and a comparison made on queueing outcomes rather than on AUROC alone. The negative result we report is about this task on these traces, where each user and each task already has enough history of its own; the only relational signal we find is on task types that have none yet.
